% Generated draft: requires mathematical and source review.
\documentclass[11pt]{article}
\usepackage[margin=25mm]{geometry}
\usepackage{fontspec}
\setmainfont{Latin Modern Roman}
\usepackage{amsmath,amssymb,mathtools,unicode-math}
\setmathfont{Latin Modern Math}
\usepackage{enumitem,needspace,xurl,hyperref}
\setlength{\parindent}{0pt}
\setlength{\parskip}{5pt}
\setlength{\emergencystretch}{3em}
\newcommand{\N}{\mathbb{N}}
\newcommand{\Z}{\mathbb{Z}}
\newcommand{\Q}{\mathbb{Q}}
\newcommand{\R}{\mathbb{R}}
\newcommand{\C}{\mathbb{C}}
\newcommand{\F}{\mathbb{F}}
\newcommand{\A}{\mathbb{A}}
\newcommand{\PP}{\mathbb{P}}
\newcommand{\E}{\mathbb{E}}
\newcommand{\eps}{\varepsilon}
\newcommand{\dd}{\,\mathrm{d}}
\begin{document}
\section*{UnsolvedMath ID: 20000335}
\textbf{Source problem code:} catalogue entry
\subsection{Definitions and mathematical statement}
The original asks for applications of stable-surface moduli to surface theory, analogous to using curve degenerations for Brill--Noether statements. It specifies no single theorem or class of applications. A stable surface conventionally has semi-log-canonical singularities and ample $\Q$-Cartier canonical divisor; the precise moduli and field setup are omitted.

The supplied report proposes a complex double-plane example, so that the topological signature and characteristic-zero double-cover interpretation are meaningful. For each integer $d\geq4$, take a smooth plane curve $C=\{f_d=0\}$ and an appropriate general degree-$2d$ polynomial $g_{2d}$. It describes the family germ
\[
 z^2=f_d^2+t g_{2d},
\]
with $z$ of weight $d$, as flat, Gorenstein, and stable, with reduced nonnormal central fiber $X_d=\PP^2\cup_C\PP^2$ and smooth double-plane general fibers. The precise coefficient and parameter-open hypotheses are not supplied; this is not asserted for every fiber over an unrestricted parameter line.

For the smooth fibers it reports
\[
\begin{gathered}
 K^2=2(d-3)^2,\qquad q=0,\qquad
 p_g=(d-1)(d-2)/2=g(C),\\
 \chi(\mathcal O)=(d^2-3d+4)/2,\qquad
 c_2=4d^2-6d+6,\qquad\operatorname{signature}=2-2d^2,
\end{gathered}
\]
with Noether defect $K^2-2p_g+4=(d-4)(d-5)$. Here $q=h^1(\mathcal O)$ and $p_g=h^0(\omega)$.

Put $V_j=H^0(\PP^2,\mathcal O(j))$ and let $T^1_{X_d}$ denote the local smoothing sheaf, conventionally $\mathcal E xt^1(\Omega^1_{X_d},\mathcal O_{X_d})$. The report identifies the normal quotient to the square-branch locus at $[f_d^2]$ with
\[
 V_{2d}/f_dV_d\cong H^0(C,T^1_{X_d}),\qquad
 \dim H^0(C,T^1_{X_d})=3d(d+1)/2.
\]
Its assertion of canonicity needs the chosen equation or the usual projective tangent-line factor made explicit. It claims every global local-smoothing section lifts within the branch chart, and a generic normal direction has $2d^2$ simple zeros producing exactly $2d^2$ ordinary double points of the total threefold. Positive degree of $T^1$ is reported to obstruct a smooth total space while this reduced central model is fixed; it does not obstruct the claimed smoothing with singular total space or variants changing the central model. These form an unreviewed boundary-to-surface example rather than a complete solution of the broad application program.
\subsection{Short English statement}
Use stable-surface boundary degenerations to obtain surface-theoretic information. The report proposes split double planes, computes their smooth-fiber invariants, and relates branch directions to local smoothing sections and total-space nodes. It distinguishes smoothing the fibers from having a smooth total threefold.
\subsection{Sources}
\begin{itemize}
\item Local UnsolvedMath registry, permanent ID $20000335$, source code \texttt{catalogue entry}; created and updated 2026-08-14. Status \texttt{partially\_solved}; classification \texttt{PARTIAL-PROGRESS}; difficulty L3; no proposer or year supplied.
\item AIM Problem Lists, Algebraic geometry, workshop \emph{Compact moduli spaces and birational geometry}, Problem 4.1, source type PDF, \url{https://aimath.org/WWN/birational/birational.pdf}. Canonical location \texttt{aim-algebraic-geometry-notes.json}, \texttt{notes[334]}, tag \texttt{problem}; statement status \texttt{exact}; original, clean, and displayed text agree. Original remarks and literature fields empty.
\item Supplied research attempt 1: result \texttt{partial\_result}, run status \texttt{partially\_solved}, rigor label \texttt{proved}, low novelty confidence. Candidate contribution \texttt{normal\_space\_smoothing\_lemma}: branch quotient, all-section lifting, normal-space dimension, $2d^2$ node count, and Noether defect. No exact primary citations or full family construction are supplied. Explicit caveat: machine-generated research aid, not peer reviewed; independently verify before citation.
\item Dataset attribution: Ulam AI / UnsolvedMath Contributors, CC BY 4.0; original problem attribution remains to the supplied AIM workshop collection.
\end{itemize}
\end{document}
